\documentclass[10pt]{beamer}
\mode<beamer>{%
  \usetheme[]{CambridgeUS}}
\usepackage{lmodern, amssymb,amsmath, graphicx, hyperref}
\usepackage[normalem]{ulem}
\usepackage{csquotes}

 \title{How Legislators Actually Invest Their Time:} 
 \subtitle{Constituents \texorpdfstring{\sout{or}}{or}\textcolor{black}{and} Policy?}
    \author{Devin Judge-Lord\thanks{Ph.D. Candidate, University of Wisconsin-Madison}\and Justin Grimmer\thanks{Associate Professor, Stanford University}, \and Eleanor Neff Powell\thanks{Booth Fowler Associate Professor, University of Wisconsin-Madison}}
    \date{\today}
    
\graphicspath{{gh-pages/Figs/}}
\begin{document}


\begin{frame}
\maketitle
\end{frame}

% Hi, My name is Devin Judge-Lord and I am presenting work with coauthors Justin Grimmer and Ellie Powell based on a very large new database of requests sent from individual member of Congress to directly to government agencies. We have spent the past 2 years collecting and coding these data, and I am excited to share our very first findings from this project. 

% We started with an a common intuition that Members of Congress have limited time and attention to divide among, for example, consgtituent work and policy work. We have come to belevie that in important ways, this is not correct. 
% Not only are legislators who do more constituntly service more likely to achive power in Washingtion, once they achieve power, they use it to advocate both policy agendas AND individual constituents.

% I am going to introduce our data and show how, in fact, the same members that do more policy work also do more constitituent service and that both seem to be connected to institutional power. 

% But first, I want to talk a little bit about our definitions and intuitions constituency service.

% From where did our incorrect intuition arise? [SLIDE]
\section{} 
\begin{frame} 
\frametitle{A constituent-policy trade-off?}
\begin{itemize}
\item ``Potomac fever'' (Fenno, 1978) \pause
% this intuition emerges from scholarship. Fenno it Potomac Feaver--the idea that the deeper a legislator gets into Washington policymakig, the less time they had for their district consticuency. 
\item ``Insiders'' who have lost touch \pause
% we also see it in popular press and campaign rhetoric contrasting washington insiders and policy wonks with outsiders who are more in touch with their constituents. 
\item \#DrainTheSwamp
% It is evoked in rallying cries to `drain the swamp' of a policy-focused elite who no longer serve their consituents.
\end{itemize}
\end{frame}


% Surely, so of this rhetoric is accusing legislators of not voting how their constituents would want. But voting is not the only way legislators represent. 
% Indeed there are far more intimate forms of representation that are less constrained by parties and chamber proceedure.
% Possibly the most common act of representation is helping constituents deal with the bureaucracy. This has been true since legislators helped constituents get pensions after the revolutionary war and especially with the rise of the administrative state. Federal agency shape nearly every aspect of our lives: 

% Thus, we focues on the requests legislators and their staff send to federal agencies.
% This aspect of representation has only begun to be studied, but has already yileded 

% [SLIDE} Why lobby the bureaucracy? 

\section{Lobbying the bureaucracy}
\begin{frame}
\frametitle{Why lobby the bureaucracy?}
Exciting new research: 
% Data on legislator requests to agencies has opened a whole new way to study representation: 
\begin{itemize}
\item Represent constituents' policy demands despite party pressure (Ritchie, 2017) \pause
% Ritchie, focusing on policy issues, finds that legislators directly engage the bureaucracy on policy issues where they face cross pressures in order to represent constituts when the party prohibits them from doing so legislativly. 
\item Constituents and institutional position, not ideology, drives policy oversight (Lowande, 2018) \pause
% Lowande, also focusing on policy finds that partisan disagreement the executive branch does not drive oversight. 
\item Discriptive = substantive (Lowande, Ritchie, Lauterbach, 2018)
% Lowande, Ritchie, and Lauterbach show that legislators make more requests on behalf of constituents whith whom they share marked identities

\item Bureaucrats have reasons to comply. 
  \begin{itemize} 
    \item And they do (Ritchie and You, n.d.). 
    \item Though not always (Mills and Kalaf-Hughes, 2015))
  \end{itemize}
\end{itemize}
\end{frame}
% Recent work by Ritchie and You and Mills and Kalaf-Hughes as has begun to investigate whether legislator requests matter and find that they do, an


% Building on this work, we focus squarely on constituent service [SLIDE]


\begin{frame}
\frametitle{What is Constituent Service?}


\begin{displayquote}
``help to individuals, groups, and localities in coping with the federal government.'' (Fenno, 1978: p. 101)
\end{displayquote} \pause
% We take a broad definition where we focus less on defining who is a constituent and more on the act of service. By this broad notion, it may be helping individuals, organized groups like unions, unorganized groups like hurricane affected populations, or companies. the key is that it is particuarlistic: an instance of someone needing help with an immigration case, grant application, or contract:
\underline{Service} 
\begin{itemize}
\item Individual casework
\item Organized groups
\item Unorganized groups
\item Companies
\end{itemize}
\end{frame}


% To study who does constituent service and when they do it, we leverage a new database of over 150,000 requests--letters, emails, phone calls that members of congress have made to over 60 agencies accross nearly all parts of the federal bureacracy.


\section{Data} 
\begin{frame}
\frametitle{Correspondence Data}



\end{frame}


% [SLIDE]


\begin{frame}
\frametitle{Correspondence Data}

\begin{itemize}
\item >250 requests to 234 federal agencies (a census)
\item >150,000 legislator requests (and counting)
\end{itemize}

\tiny

\begin{table}[!htbp] \centering 
  \caption{Responses to FOIA Requests for a Census of Congressional Contacts Thus Far} 
  \label{responserates}
\begin{tabular}{@{\extracolsep{5pt}} lccc} 
%\\[-1.8ex]\hline 
%\hline \\[-1.8ex] 
Department & FOIAed Components & Provided Records & Coded \\ 
\hline \\[-1.8ex] 
Department of Agriculture & 29 & 29 & 11 \\ 
Department of Commerce & 18 & 18 & 10 \\ 
Department of Defense & 49 & 10 & 7 \\ 
Department of Education & 1 & 1 & 1 \\ 
Department of Energy & 8 & 2 & 1 \\ 
Department of Health and Human Services & 15 & 7 & 5 \\ 
Department of Homeland Security & 5 & 4 & 3 \\ 
Department of Housing and Urban Development & 2 & 0 & 0 \\ 
Department of Justice & 11 & 1 & 1 \\ 
Department of Labor & 24 & 9 & 6 \\ 
Department of State & 1 & 0 & 0 \\ 
Department of the Interior & 10 & 6 & 5 \\ 
Department of the Treasury & 8 & 5 & 4 \\ 
Department of Transportation & 10 & 6 & 6 \\ 
Department of Veterans Affairs & 3 & 0 & 0 \\ 
Independent Agencies & 34 & 21 & 18 \\ 
\hline \\[-1.8ex] 
\end{tabular} 
\end{table} 

\end{frame}


\section{Results: Institutional power and constituency service}
\begin{frame}
\frametitle{Institutional Power and Constituency Service}
% So this is the part where we discover that we were wrong. 
% We found we found the same powerful members of Congress who are most acively involved in policy oversight are also those most engaged in helping constituents. 
% First, lets look at a simple cross sectional comparisons of members who do and do 
Legislators who have power advocate more on behalf of constituents.

\includegraphics[width = \textwidth]{cross_sectional_models-1}
\end{frame}




\begin{frame}
\frametitle{Difference-in-Difference in Constituency Service}

Legislators who \textcolor{blue}{gain} power advocate more on behalf of constituents.

\includegraphics[width = \textwidth]{DiD-1}
\end{frame}


\begin{frame}
\frametitle{Statutory Oversight Power}

\begin{itemize}
\item Agencies receive more constituent service requests from chairs of their oversight committees
\item Legislators who become chairs then lobby more for constituents, targeting the agencies they oversee
\end{itemize}

\includegraphics[width = \textwidth]{oversight_committee_chair_models-1}

\end{frame}


\end{document}